% Minimal preamble for Open Logic Project sources (pinned commit 1e960beff9ed7835bf3e3f1335e21af3439cd107).
%
% OLP's own master preamble cannot run on the project host: TeX Live 2025 there lacks
% prooftree.sty, ntheorem.sty, extarrows.sty, mathdots.sty, isomath.sty and subfiles.sty.
% Following the project's own rule ("if the author preamble needs packages the toolchain lacks,
% do not invent macros -- use a minimal working preamble"), this file declares, verbatim in
% spirit and numbering behaviour, exactly the theorem-like environments OLP defines in
% sty/open-logic-book-envs.sty (thm/lem/prop/cor/defn/ex/prob/rem/note/case/axiom/conv),
% plus the shorthand accents OLP uses (`!!a`, `!!b`, `!!{...}`).
%
% Author text is untouched; only the packaging is ours.
\documentclass[10pt,a4paper]{article}

\usepackage{amsmath}
\usepackage{amssymb}
\usepackage{amsthm}
\usepackage{graphicx}
\usepackage{hyperref}
\usepackage{enumitem}
\usepackage[framemethod=default]{mdframed}

% ---- OLP shorthand accents (from sty/open-logic.sty) ----
% NB: \! is taken by LaTeX (negative thin space), so we only define the !!a / !!b families
% and the \olup helper used by them.
\newcommand*{\olup}[1]{\MakeUppercase{#1}}
\def\!!!a#1{\olup{#1}}
\def\!!a#1{\olup{#1}}
\def\!!b#1{\olup{#1}}

% ---- theorem-like environments, same family numbering as OLP ----
\theoremstyle{plain}
\newtheorem{thm}{Theorem}[section]
\newtheorem{lem}[thm]{Lemma}
\newtheorem{prop}[thm]{Proposition}
\newtheorem{cor}[thm]{Corollary}
\newtheorem{axiom}{Axiom}

\theoremstyle{definition}
\newtheorem{defn}[thm]{Definition}
\newtheorem{ex}[thm]{Example}
\newtheorem{prob}{Problem}[section]
\newtheorem{probd}{Problem}
\newtheorem{conv}[thm]{Convention}

\theoremstyle{remark}
\newtheorem{rem}[thm]{Remark}
\newtheorem{note}[thm]{Note}
\newtheorem{case}[thm]{Case}

% ---- environments OLP defines as wrappers ----
\newenvironment{probtag}[1]{\begin{prob}}{\end{prob}}
\newenvironment{explain}{}{}
\newenvironment{explain*}{}{}
\newenvironment{editorial}{}{}
\newenvironment{digress}{}{}
\newenvironment{digress*}{}{}
\newenvironment{history}{}{}
\newenvironment{reference}{}{}
\newcommand*{\olfileid}[3]{}
\newcommand*{\olsection}[1]{\section*{#1}}
\newcommand*{\olchapter}[3]{}
\newcommand*{\olpart}[2]{}
\newcommand*{\olimport}[1]{}
\newcommand*{\olinclude}[1]{}
\newcommand*{\olpath}[1]{}
\newcommand*{\ollangid}[1]{}
\newcommand*{\problemtag}[1]{}
\newcommand*{\problemauthor}[1]{}
\newcommand*{\ollabel}[1]{\label{#1}}
% \olref takes 0-3 OPTIONAL arguments plus one mandatory label:
%     \olref{lab}  |  \olref[from]{lab}  |  \olref[from][to]{lab}  |  \olref[a][b][c]{lab}
% A single-argument definition silently swallows "[sth]" and emits a broken reference
% (measured on the host: the text came out as "??iso]x"). The chain below consumes the
% brackets and cross-references only the mandatory argument.
\makeatletter
\newcommand{\olref}{\@ifnextchar[{\olref@i}{\olref@final}}
\def\olref@i[#1]{\@ifnextchar[{\olref@ii}{\olref@final}}
\def\olref@ii[#1]{\@ifnextchar[{\olref@iii}{\olref@final}}
\def\olref@iii[#1]{\olref@final}
\def\olref@final#1{\ref{#1}}
\makeatother
\newcommand*{\olnewpage}{\clearpage}

% ===================== OLP compatibility layer (1) =====================
% Compatibility layer for Open Logic Project sources (pinned commit 1e960beff9ed7835bf3e3f1335e21af3439cd107).
%
% OLP's own style stack (sty/open-logic.sty) needs ntheorem, prooftree, extarrows, mathdots,
% isomath and subfiles; the project's local and host TeX Live installs lack several of those.
% Rather than vendor packages, this file supplies drop-in equivalents for every OLP construct
% that the frozen units actually use. It only rearranges the author's own tokens; no
% mathematical content is added, and nothing here is claimed to be OLP's own file.

\makeatletter

% ---------------- proof trees (OLP's prooftree is a bussproofs derivative) ----------------
% These commands CONSUME their argument and typeset nothing: the bracket labels of OLP's proof
% trees carry rule names. Printing them inline dropped those labels into math mode and produced
% 63 "Missing $" / "allowed only in math mode" errors; swallowing them keeps the brace balance
% and the mode correct. The derivation text between the commands is still printed unchanged.
\providecommand{\AxiomC}[1]{}
\providecommand{\UnaryInfC}[1]{}
\providecommand{\BinaryInfC}[1]{}
\providecommand{\TrinaryInfC}[1]{}
\providecommand{\QuaternaryInfC}[1]{}
\providecommand{\QuinaryInfC}[1]{}
\providecommand{\RightLabel}[1]{}
\providecommand{\LeftLabel}[1]{}
\providecommand{\doubleLine}{}
\providecommand{\noLine}{}
\providecommand{\DisplayProof}{}
\providecommand{\fCenter}{}
\newenvironment{prooftree}{}{}
\newenvironment{ltprooftree}{}{}
\newenvironment{oltableau}{}{}

% ---------------- logical notation ----------------
\providecommand{\lif}{\mathbin{\rightarrow}}
\providecommand{\lforall}{\mathord{\forall}}
\providecommand{\lexists}{\mathord{\exists}}
\providecommand{\ltrue}{\top}
\providecommand{\lfalse}{\bot}
\providecommand{\Ltrue}{\top}
\providecommand{\Lfalse}{\bot}
\providecommand{\Sequent}{\mathrel{\Rightarrow}}
\providecommand{\Proves}{\mathrel{\vdash}}
\providecommand{\Entails}{\mathrel{\models}}
\providecommand{\Subst}[3]{#1[#2/#3]}
\providecommand{\sFmla}[2]{#1(#2)}
\providecommand{\pFmla}[2]{#1(#2)}
\providecommand{\formula}{\relax}
\providecommand{\True}{\top}
\providecommand{\False}{\bot}
\providecommand{\Args}[1]{#1}
\providecommand{\Ts}{\ensuremath{\top}}
\providecommand{\Fi}{\ensuremath{\bot}}
\providecommand{\lfalsum}{\bot}

% ---------------- objects, structures, metatheory ----------------
\providecommand{\Obj}{}
\providecommand{\Struct}{\ensuremath{\mathfrak{A}}}
\providecommand{\Domain}{\ensuremath{|\mathfrak{A}|}}
\providecommand{\Assign}{\ensuremath{s}}
\providecommand{\Sat}[2]{\ensuremath{#1 \models #2}}
\providecommand{\mSat}[3]{\ensuremath{#1 \models #2[#3]}}
\providecommand{\mModel}[2]{\ensuremath{#1 \models #2}}
\providecommand{\Lang}{\ensuremath{\mathcal{L}}}
\providecommand{\Th}{\ensuremath{\mathrm{Th}}}
\providecommand{\Log}{\ensuremath{\mathrm{L}}}
\providecommand{\tuple}[1]{\langle #1 \rangle}
\providecommand{\num}[1]{\underline{#1}}
\providecommand{\cfind}[1]{\varphi_{#1}}
\providecommand{\Nat}{\ensuremath{\mathbb{N}}}
\providecommand{\Mult}{\ensuremath{\mathrm{Mult}}}
\providecommand{\Add}{\ensuremath{\mathrm{Add}}}
\providecommand{\Exp}{\ensuremath{\mathrm{Exp}}}
\providecommand{\Zero}{\ensuremath{0}}
\providecommand{\Succ}{\ensuremath{\mathrm{S}}}
\providecommand{\lambd}{\ensuremath{\lambda}}
\providecommand{\bredpar}{\mathrel{\rightarrow_\beta}}
\providecommand{\bcd}[1]{\ensuremath{\ulcorner #1 \urcorner}}
\providecommand{\Set}[1]{\ensuremath{\{#1\}}}
\providecommand{\indcase}[1]{\par\noindent\emph{Case #1}.\quad}
\providecommand{\indfrm}[1]{\par\noindent\emph{Form #1}.\quad}
\providecommand{\iftag}[1]{#1}
\providecommand{\tagitem}[1]{\item[#1]}
% these five are used with arguments by OLP's proof trees -> swallow the argument
\providecommand{\Elim}[1]{}
\providecommand{\Intro}[1]{}
\providecommand{\LeftR}[1]{}
\providecommand{\RightR}[1]{}
\providecommand{\TAss}[1]{}
\providecommand{\TRule}[1]{}

\makeatother
\providecommand{\olpath}{.}

% ===================== OLP compatibility layer (2) =====================
% --- second compatibility block: notation and shorthands used by the frozen OLP units ---
% Everything here is \providecommand, so anything already defined by LaTeX or the packages wins.
\providecommand{\CutCS}{\mathrm{Cut}}
\providecommand{\LeftSubproofLabel}{\relax}
\providecommand{\RightSubproofLabel}{\relax}
\providecommand{\maeh}{\relax}
\providecommand{\cheight}{\relax}
\providecommand{\ande}{\mathrm{and}_E}
\providecommand{\andi}{\mathrm{and}_I}
\providecommand{\ore}{\mathrm{or}_E}
\providecommand{\ori}{\mathrm{or}_I}
\providecommand{\typeof}{\mathord{:}}
\providecommand{\pValue}[1]{#1}
\providecommand{\Sent}{\mathrm{Sent}}
\providecommand{\cardplus}{\mathbin{+}}
\providecommand{\cardtimes}{\mathbin{\cdot}}
\providecommand{\cardexpo}{\mathbin{\hat{\ }}}
\providecommand{\cardnless}{\not<}
\providecommand{\cardle}{\leq}
\providecommand{\funfromto}[3]{#1\colon #2 \to #3}
\providecommand{\funimage}[1]{\operatorname{im}(#1)}
\providecommand{\fdefined}{\mathord{\downarrow}}
\providecommand{\canonord}{\operatorname{ord}}
\providecommand{\ordtimes}{\mathbin{\cdot}}
\providecommand{\ordeq}[2]{#1 = #2}
\providecommand{\ordneq}[2]{#1 \neq #2}
\providecommand{\ordtype}[1]{\operatorname{ord}(#1)}
\providecommand{\ordsucc}{\mathord{+1}}
\providecommand{\ordunion}{\cup}
\providecommand{\onesphere}{S^1}
\providecommand{\rotationsgroup}{\mathrm{SO}(2)}
\providecommand{\equivrep}[1]{[#1]}
\providecommand{\rlexless}{\prec}
\providecommand{\ZFC}{\mathrm{ZFC}}
\providecommand{\ZF}{\mathrm{ZF}}
\providecommand{\Rat}{\mathbb{Q}}
\providecommand{\Real}{\mathbb{R}}
\providecommand{\PosInt}{\mathbb{Z}^{+}}
\providecommand{\Id}{\mathrm{id}}
\providecommand{\Bin}{\mathrm{Bin}}
\providecommand{\pto}{\rightharpoonup}
\providecommand{\dom}{\operatorname{dom}}
\providecommand{\ran}{\operatorname{ran}}
\providecommand{\len}{\operatorname{len}}
\providecommand{\concat}{\mathbin{{}^\frown}}
\providecommand{\card}[1]{\lvert #1 \rvert}
\providecommand{\Pow}{\mathcal{P}}
\providecommand{\Setabs}[2]{\{\, #1 \mid #2 \,\}}
\providecommand{\comp}{\circ}
\providecommand{\Char}{\chi}
\providecommand{\Prf}{\operatorname{Prf}}
\providecommand{\Prov}{\operatorname{Prov}}
\providecommand{\Proj}{\operatorname{proj}}
\providecommand{\Gn}{\operatorname{gn}}
\providecommand{\gn}{\operatorname{gn}}
\providecommand{\Value}[1]{\operatorname{val}(#1)}
\providecommand{\Deduce}{\mathrm{D}}
\providecommand{\Dual}{\mathrm{d}}
\providecommand{\Ax}{\mathrm{Ax}}
\providecommand{\pSat}[3]{#1 \models #2\,[#3]}
\providecommand{\pAssign}{\sigma}
\providecommand{\bforall}{\forall}
\providecommand{\bexists}{\exists}
\providecommand{\liff}{\mathbin{\leftrightarrow}}
\providecommand{\FV}{\operatorname{FV}}
\providecommand{\UnaryInf}{\ensuremath{\mid}}
\providecommand{\BinaryInf}{\ensuremath{\mid}}
\providecommand{\depth}{\operatorname{depth}}
\providecommand{\oliflabeldef}[3]{#1}
\providecommand{\TMblank}{\sqcup}
\providecommand{\TMstroke}{\mid}
\providecommand{\TMstay}{\mathrm{S}}
\providecommand{\TMright}{\mathrm{R}}
\providecommand{\TMleft}{\mathrm{L}}
\providecommand{\TMendtape}{\triangleright}
\providecommand{\umin}{\mu}
\providecommand{\Sum}{\sum}
% bibliography helpers (natbib): the frozen units cite the shared OLP bibliography
\providecommand{\citep}[1]{\cite{#1}}
\providecommand{\citet}[1]{\cite{#1}}
\providecommand{\citeauthor}[1]{\cite{#1}}
\providecommand{\citeyear}[1]{\cite{#1}}

% ===================== OLP compatibility layer (3) =====================
% --- third compatibility block: remaining notation seen in the frozen OLP units ---
\providecommand{\iso}{\cong}
\providecommand{\nszero}{\mathrm{0}}
\providecommand{\nssucc}{\mathrm{S}}
\providecommand{\nsplus}{\mathbin{+}}
\providecommand{\nssless}{<}
\providecommand{\nsless}{<}
\providecommand{\PL}{\mathrm{PL}}
\providecommand{\Nec}{\mathord{\Box}}
\providecommand{\RK}{\mathrm{RK}}
\providecommand{\ST}{\mathrm{ST}}
\providecommand{\SSubst}[3]{#1[#2/#3]}
\providecommand{\subst}[3]{#1[#2/#3]}
\providecommand{\tagfalse}{\mathrm{F}}
\providecommand{\cutr}{\mathrm{Cut}}
\providecommand{\cardneq}{\neq}
\providecommand{\issucc}[1]{\operatorname{succ}(#1)}
\providecommand{\tup}[1]{\langle #1 \rangle}
\providecommand{\Var}{\mathrm{Var}}
\providecommand{\Bew}{\operatorname{Bew}}
\providecommand{\Con}{\operatorname{Con}}
\providecommand{\Cons}{\operatorname{Cons}}
\providecommand{\pDeduce}{\mathrm{D}}
\providecommand{\psat}{\models}
\providecommand{\den}[1]{\llbracket #1 \rrbracket}
\providecommand{\Sem}[2]{\llbracket #1 \rrbracket_{#2}}
\providecommand{\DomainOf}[1]{\lvert #1 \rvert}
\providecommand{\Form}{\mathrm{Form}}
\providecommand{\Terms}{\mathrm{Term}}
\providecommand{\Atom}{\mathrm{Atom}}
\providecommand{\FVars}[1]{\operatorname{FV}(#1)}
\providecommand{\fv}[1]{\operatorname{fv}(#1)}
\providecommand{\tuplep}[1]{\langle #1 \rangle}
\providecommand{\restrict}[2]{#1 \upharpoonright #2}
\providecommand{\sto}{\rightarrow}
\providecommand{\FmlaSet}[1]{\mathrm{Fmla}_{#1}}
\providecommand{\Thm}{\mathrm{Thm}}
\providecommand{\Eval}{\operatorname{val}}
\providecommand{\blank}{\sqcup}
\providecommand{\blankTape}{\sqcup}
\providecommand{\Consis}{\operatorname{Con}}
\providecommand{\drs}{\mathrm{d}}
\providecommand{\VRule}{\mid}
\providecommand{\Bet}{\mathrm{B}}
\providecommand{\Seg}[1]{\mathrm{Seg}_{#1}}
\providecommand{\Enums}{\mathrm{En}}
\providecommand{\Z}{\mathbb{Z}}
\providecommand{\N}{\mathbb{N}}
\providecommand{\abs}[1]{\lvert #1 \rvert}
\providecommand{\floor}[1]{\lfloor #1 \rfloor}
\providecommand{\ceil}[1]{\lceil #1 \rceil}

% ===================== OLP compatibility layer (4) =====================
% --- fourth compatibility block: harvested from the per-candidate compile logs ---
\providecommand{\beredpar}{\mathrel{\rightarrow_\beta}}
% \bered / \bred are used as infix relations AND with a bracketed annotation, so they need the
% optional-argument chain; \providecommand would leave an argument-less definition in place.
\makeatletter
\def\bered{\@ifnextchar[{\bered@i}{\mathrel{\rightarrow_\beta}}}
\def\bered@i[#1]{\mathrel{\rightarrow_\beta}}
\def\bred{\@ifnextchar[{\bred@i}{\mathrel{\rightarrow_\beta}}}
\def\bred@i[#1]{\mathrel{\rightarrow_\beta}}
\makeatother
\providecommand{\beredone}{\mathrel{=_\beta}}
\providecommand{\bredone}{\mathrel{=_\beta}}
\providecommand{\xredone}{\mathrel{=_\beta}}
\providecommand{\aconvone}{\mathrel{=_\alpha}}
\providecommand{\cardless}{<}
\providecommand{\tagrefs}[1]{#1}
\providecommand{\ordplus}{\mathbin{+}}
\providecommand{\funrestrictionto}[2]{#1 \upharpoonright #2}
\providecommand{\fundefined}{\mathord{\uparrow}}
\providecommand{\Trm}{\mathrm{Trm}}
\providecommand{\OProv}{\operatorname{Prov}}
\providecommand{\OPrf}{\operatorname{Prf}}
\providecommand{\aeq}{\mathrel{=}}
\providecommand{\mClass}[1]{[#1]}
\providecommand{\equivclass}[1]{[#1]}
\providecommand{\pheight}{\operatorname{ht}}
\providecommand{\phy}{\operatorname{ht}}
\providecommand{\Cut}{\mathrm{Cut}}
\providecommand{\QuantRank}{\operatorname{qr}}
\providecommand{\elemequiv}{\mathrel{\equiv}}
\providecommand{\PIso}{\mathrm{PIso}}
\providecommand{\Theory}{\mathrm{Th}}
\providecommand{\supstrict}{\subsetneq}
\providecommand{\trcl}{\operatorname{trcl}}
\providecommand{\cref}[1]{\ref{#1}}
\providecommand{\Discharge}{\mathrm{Dis}}
\providecommand{\Weakening}{\mathrm{Wk}}
\providecommand{\ZFminus}{\mathrm{ZF}^{-}}
\providecommand{\RK}{\mathrm{RK}}
% \Dual is printed both in math mode and (with a bracketed annotation) in running text;
% the math-mode test keeps \mathrm legal in both cases.
\makeatletter
\renewcommand{\Dual}{\@ifnextchar[{\Dual@i}{\ifmmode\mathrm{d}\else$\mathrm{d}$\fi}}
\def\Dual@i[#1]{\ifmmode\mathrm{d}[#1]\else$\mathrm{d}[#1]$\fi}
\makeatother
\providecommand{\RightR}{\mathrm{R}}
\providecommand{\LeftR}{\mathrm{L}}
\providecommand{\Intro}{\mathrm{I}}
\providecommand{\Axiom}{\mathrm{Ax}}
\providecommand{\formula}[1]{#1}
\providecommand{\stared}{\ast}
\providecommand{\oset}[2]{\{\, #1 \mid #2 \,\}}
\providecommand{\setof}[2]{\{\, #1 \mid #2 \,\}}
\providecommand{\pow}[1]{\mathcal{P}(#1)}
\providecommand{\ordpair}[2]{\langle #1, #2 \rangle}
\providecommand{\ordtriple}[3]{\langle #1, #2, #3 \rangle}
\providecommand{\union}{\cup}
\providecommand{\inter}{\cap}
\providecommand{\seq}[1]{\langle #1 \rangle}
\providecommand{\SuchThat}{\mid}
\providecommand{\where}{\mid}
\providecommand{\TVar}{\mathrm{TVar}}
\providecommand{\Bool}{\mathrm{Bool}}
\providecommand{\Val}{\operatorname{Val}}
\providecommand{\FV}{\operatorname{FV}}
\providecommand{\TL}[1]{\mathrm{TL}}
\providecommand{\xstared}{\ast}

% ===================== OLP compatibility layer (5) =====================
% --- fifth compatibility block: rescues the quarantined OLP candidates ---
% The proof-tree and mode fixes live in block 1 (where those macros are first defined); this block
% only adds notation that the quarantined spans still needed and that no earlier block covered.
\providecommand{\ident}{\operatorname{id}}
\providecommand{\setrank}{\operatorname{rank}}
\providecommand{\cardfont}{\mathfrak{c}}
\providecommand{\Complement}{\overline}
\providecommand{\QR}{\mathrm{QR}}
\providecommand{\becd}{\ensuremath{\ulcorner}}
\providecommand{\disjointsum}{\mathbin{\oplus}}
\providecommand{\cardeq}[2]{#1 = #2}
\providecommand{\xred}{\mathrel{\rightarrow}}
\providecommand{\cardless}{<}
\providecommand{\cardplus}{\mathbin{+}}
\providecommand{\cardtimes}{\mathbin{\cdot}}
\providecommand{\cardexpo}{\mathbin{\hat{\ }}}
\providecommand{\cardneq}{\neq}
\providecommand{\cardle}{\leq}
\providecommand{\cardnless}{\not<}
\begin{document}

% ==== OLP file: content/set-theory/choice/banach.tex ====


\olfileid{sth}{choice}{banach}
\olsection{The Banach--Tarski Paradox}

We might also attempt to justify Choice, as Boolos attempted to
justify Replacement, by appealing to \emph{extrinsic} considerations
(see \olref[replacement][extrinsic]{sec}). After all, adopting Choice
has many desirable consequences: the ability to compare every
cardinal; the ability to well-order every set; the ability to treat
cardinals as a particular kind of ordinal; etc. 

Sometimes, however, it is claimed that Choice has \emph{undesirable}
consequences. Mostly, this is due to a result by
\cite{BanachTarski1924}. 

\begin{thm}[Banach--Tarski Paradox (in $\ZFC$)]
Any ball can be decomposed into finitely many pieces, which can be
reassembled (by rotation and transportation) to form two copies of
that ball.
\end{thm}
\noindent 
At first glance, this is a bit amazing. Clearly the two balls have
\emph{twice} the volume of the original ball. But rigid
motions---rotation and transportation---do not change volume. So it
looks as if Banach--Tarski allows us to magick new matter into
existence.

It gets worse.\footnote{See \citet[Theorem 3.12]{Wagon2016}.} Similar
reasoning shows that a pea can be cut into finitely many pieces, which
can then be reassembled (by rotation and transportation) to form an
entity the shape and size of Big Ben.

None of this, however, holds in $\ZF$ on its own.\footnote{Though
Banach--Tarski can be proved with principles which are strictly weaker
than Choice; see \citet[303]{Wagon2016}.} So we face a decision:
reject Choice, or learn to live with the ``paradox''. 

We're going to suggest that we should learn to live with the
``paradox''. Indeed, we don't think it's much of a paradox at all. In
particular, we don't see why it is any more or less paradoxical than
any of the following results:\footnote{\citet[276--7]{Potter2004},
\citet[16]{Weston2003}, \citet[31, 308--9]{Wagon2016}, make
similar points, using other examples.}
\begin{enumerate}
	\item There are as many points in the interval $(0,1)$ as in  $\Real$. 
	\\\emph{Proof}: consider $\tan(\pi(r-\nicefrac{1}{2})))$.
	\item There are as many points in a line as in a square.
	\\See \olref[his][set][pathology]{sec} and \olref[his][set][cantorplane]{sec}.
	\item There are space-filling curves. 
	\\See \olref[his][set][pathology]{sec} and \olref[his][set][hilbertcurve]{sec}.
\end{enumerate}
None of these three results require Choice. Indeed, we now just regard
them as surprising, lovely, bits of mathematics. Maybe we should adopt
the same attitude to the Banach--Tarski Paradox.

To be sure, a technical observation is required here; but it only
requires keeping a level head. Rigid motions preserve volume.
Consequently, the five\footnote{We stated the Paradox in terms of
``finitely many pieces''. In fact, \citet{Robinson1947} proved that
the decomposition can be achieved with \emph{five} pieces
(but no fewer). For a proof, see \citet[pp.~66--7]{Wagon2016}.} pieces
into which the ball is decomposed cannot all be \emph{measurable}.
Roughly put, then, it makes no sense to assign a volume to these
individual pieces. You should think of these as unpicturable,
``infinite scatterings'' of points. Now, maybe it is ``weird'' to
conceive of such ``infinitely scattered'' sets. But their existence
seems to fall out from the injunction, embodied in \stagesacc{}, that
you should form \emph{all possible} collections of earlier-available
sets. 

If none of that convinces, here is a final (extrinsic) argument in
favour of embracing the Banach--Tarski Paradox. It immediately entails
the best math joke of all time:
\begin{enumerate}
	\item[] \emph{Question}. What's an anagram of ``Banach--Tarski''? 
	\item[] \emph{Answer}. ``Banach--Tarski Banach--Tarski''.
\end{enumerate}

% ==== OLP file: content/set-theory/choice/choice.tex ====
% Part: set-theory
% Chapter: choice



\olchapter{sth}{choice}{Choice}

\olimport{introduction}
\olimport{tarskiscott}
\olimport{hartogs}
\olimport{wellorderingproblem}
\olimport{countablechoice}
\olimport{justifications}
\olimport{banach}
\olimport{vitali}

\OLEndChapterHook

% ==== OLP file: content/set-theory/choice/countablechoice.tex ====


\olfileid{sth}{choice}{countablechoice}
\olsection{Countable Choice}

It is easy to prove, without any use of Choice/Well-Ordering, that:

\begin{lem}[in $\Zminus$]
Every finite set has a choice function. 
\end{lem}

\begin{proof}
Let $a = \{b_1, \ldots, b_n\}$. Suppose for simplicity that each $b_i
\neq \emptyset$. So there are objects $c_1, \ldots, c_n$ such that
$c_1 \in b_1, \ldots, c_n \in b_n$. Now by
\olref[z][pairs]{prop:pairsconsequences}, the set $\{\langle b_1,
c_1\rangle , \ldots, \langle b_n, c_n\rangle\}$ exists; and this is a
choice function for~$a$.
\end{proof}

But matters get murkier as soon as we consider infinite sets. For
example, consider this ``minimal'' extension to the above:

\begin{defish}
\emph{Countable Choice.} Every \emph{countable} set has a choice function. 
\end{defish}

This is a special case of Choice. And it transpires that this
principle was invoked fairly frequently, without  an obvious awareness
of its use. Here are two nice examples.\footnote{Due to
\citet[\S9.4]{Potter2004} and Luca Incurvati.}

\begin{ex}
Here is a natural thought: for any set $A$, either
$\cardle{\omega}{A}$, or $\cardeq{A}{n}$ for some $n \in \omega$. This
is one way to state the intuitive idea, that every set is either
finite or infinite. Cantor, and many other mathematicians, made this
claim without proving it. Cautious as we are, we proved this in
\olref[cardinals][classing]{generalinfinitycharacter}. But
in that proof we were working in $\ZFC$, since we were assuming that
any set $A$ can be well-ordered, and hence that $\card{A}$ is
guaranteed to exist. That is: we explicitly assumed Choice.

In fact, \citet{Dedekind1888} offered his own proof of
this claim, as follows:

\begin{thm}[in $\Zminus + \text{Countable Choice}$]
For any $A$, either $\cardle{\omega}{A}$ or $\cardeq{A}{n}$ for some
$n \in \omega$.
\end{thm}

\begin{proof}
Suppose $\cardneq{A}{n}$ for all $n \in \omega$. Then in particular
for each $n < \omega$ there is subset $A_n \subseteq A$ with exactly
$\cardexpo{2}{n}$ elements. Using this sequence $A_0, A_1, A_2,
\ldots$, we define for each $n$:
\[
	B_n = A_n \setminus \bigcup_{i < n} A_i.
\]
Now note the following
\begin{align*}
	\card{\bigcup_{i < n}A_n} 
	&\leq \card{A_0} + \card{A_1} + \ldots + \card{A_{n-1}}\\
	&=1 + 2 + \ldots + 2^{n-1}\\
	& = 2^n - 1\\
	& < 2^n = \card{A_n}
\end{align*}
Hence each $B_n$ has at least one member, $c_n$. Moreover, the $B_n$s
are pairwise disjoint; so if $c_n = c_m$ then $n = m$. But every $c_n
\in A$. So the function  $f(n) = c_n$ is an injection $\omega \to A$.
\end{proof}
\noindent 
Dedekind did not flag that he had used Countable Choice. But, did
\emph{you} spot its use? Look again. (Really: \emph{look again}.)

The proof used Countable Choice twice. We used it once, to obtain
our sequence of sets $A_0$, $A_1$, $A_2$, \dots\@ We then used it
again to select our elements $c_n$ from each~$B_n$. Moreover, this use
of Choice is ineliminable. \citet[p.~138]{Cohen1966} proved that the
result fails if we have no version of Choice. That is: it is
consistent with $\ZF$ that there are sets which are
\emph{incomparable} with~$\omega$.
\end{ex}

\begin{ex} 
In \citeyear{Cantor1878}, Cantor stated that a countable union of
countable sets is countable. He did not present a proof, perhaps
indicating that he took the proof to be obvious. Now, cautious as we
are, we proved a more general version of this result in
\olref[card-arithmetic][simp]{kappaunionkappasize}.  But our proof
explicitly assumed Choice. And even the proof of the less general
result requires Countable Choice.

\begin{thm}[in $\Zminus + \text{Countable Choice}$]
If $A_n$ is countable for each $n \in \omega$, then $\bigcup_{n <
\omega} A_n$ is countable.
\end{thm}

\begin{proof}
Without loss of generality, suppose that each $A_n \neq \emptyset$. So
for each $n \in \omega$ there is !!a{surjection} $f_n \colon \omega
\to A_n$. Define $f \colon \omega \times \omega \to \bigcup_{n <
\omega} A_n$ by $f(m, n) = f_n(m)$. The result follows because $\omega
\times \omega$ is countable
(\olref[sfr][siz][zigzag]{natsquaredenumerable}) and $f$ is
!!a{surjection}.
\end{proof}
\noindent 
Did you spot the use of the Countable Choice? It is used to choose our
sequence of functions $f_0$, $f_1$, $f_2$, \dots\footnote{A similar
use of Choice occurred in
\olref[card-arithmetic][simp]{kappaunionkappasize}, when we gave the
instruction ``For each $\beta \in \cardfont{a}$, fix !!a{injection}
$f_\beta$''.} And again, the result fails in the absence of any Choice
principle. Specifically, \citet{FefermanLevy1963} proved that it is
consistent with $\ZF$ that a countable union of countable sets has
cardinality~$\beth_1$. But here is a much funnier statement of the
point, from Russell:
\begin{quote}
  This is illustrated by the millionaire who bought a pair of socks
  whenever he bought a pair of boots, and never at any other time, and
  who had such a passion for buying both that at last he had
  $\aleph_0$ pairs of boots and $\aleph_0$ pairs of socks\dots\@ Among
  boots we can distinguish right and left, and therefore we can make a
  selection of one out of each pair, namely, we can choose all the
  right boots or all the left boots; but with socks no such principle
  of selection suggests itself, and we cannot be sure, unless we
  assume the multiplicative axiom [i.e., in effect Choice], that there
  is any class consisting of one sock out of each pair.
  \citep[p.~126]{Russell1919}
\end{quote}
In short, some form of Choice is needed to prove the following: If you
have countably many pairs of socks, then you have (only) countably
many socks. And in fact, without Countable Choice (or something
equivalent), a countable union of countable sets can fail to be
countable. 
\end{ex}

The moral is that Countable Choice was used repeatedly, without much
awareness of its users. The philosophical question is: How could we
\emph{justify} Countable Choice? 

An attempt at an intuitive justification might invoke an appeal to a
supertask. Suppose we make the first choice in $\nicefrac{1}{2}$ a
minute, our second choice in $\nicefrac{1}{4}$ a minute, \dots, our
$n$-th choice in $\nicefrac{1}{2^n}$ a minute, \dots\@ Then within $1$~minute, we will have made an $\omega$-sequence of choices, and defined
a choice function. 

But what, really, could such a thought-experiment tell us? For a
start, it relies upon taking this idea of ``choosing'' rather
literally. For another, it seems to bind up mathematics in
metaphysical possibility. 

More important: it is not going to give us any justification for
Choice \emph{tout court}, rather than \emph{mere} Countable Choice.
For if we need \emph{every} set to have a choice function, then we'll
need to be able to perform a ``supertask of arbitrary ordinal
length.'' Bluntly, that idea is laughable.

% ==== OLP file: content/set-theory/choice/hartogs.tex ====


\olfileid{sth}{choice}{hartogs}
\olsection{Comparability and Hartogs' Lemma}

That's the plus side. Here's the minus side. Without Choice, things
get \emph{messy}. To see why, here is a nice result due to
\cite{Hartogs1915}:

\begin{lem}[\emph{in $\ZF$}]\ollabel{HartogsLemma}
For any set $A$, there is an ordinal $\alpha$ such that $\cardnless{\alpha}{A}$
\end{lem}

\begin{proof}
If $B \subseteq A$ and $R \subseteq B^2$, then $\tuple{B, R} \subseteq
V_{\setrank{A}+4}$ by
\olref[ord-arithmetic][using-addition]{rankcomputation}. So, using
Separation, consider:
\[
	C = \Setabs{\tuple{B, R} \in V_{\setrank{A}+5}}{B\subseteq A 
	\text{ and $\tuple{B, R}$ is a well-ordering}}
\]
Using Replacement and
\olref[ordinals][ordtype]{thmOrdinalRepresentation}, form the set: 
\[
	\alpha = \Setabs{\ordtype{B, R}}{\tuple{B, R} \in C}.
\]
By \olref[ordinals][basic]{corordtransitiveord}, $\alpha$ is an
ordinal, since it is a transitive set of ordinals. After all, if
$\gamma \in \beta \in \alpha$, then $\beta = \ordtype{B, R}$ for some
$B \subseteq R$, whereupon $\gamma = \ordtype{B_b, R_b}$ for some $b
\in B$ by \olref[ordinals][iso]{wellordinitialsegment}, so that
$\gamma \in \alpha$. 

For reductio, suppose there is !!a{injection} $f \colon \alpha \to A$.
Then, where:
\begin{align*}
	B &= \ran{f}\\
	R &= \Setabs{\tuple{f(\alpha), f(\beta)} \in A \times A}{\alpha \in \beta}.
\end{align*}
Clearly $\alpha = \ordtype{B, R}$ and $\tuple{B, R} \in C$. So $\alpha
\in \alpha$, which is a contradiction.
\end{proof}

This entails a deep result:

\begin{thm}[\emph{in $\ZF$}]
The following claims are equivalent:
\begin{enumerate}
	\item\ollabel{equivwo} The Axiom of Well-Ordering
	\item\ollabel{equivcompare} Either $\cardle{A}{B}$ or
	$\cardle{B}{A}$, for any sets $A$ and $B$
\end{enumerate}
\end{thm}

\begin{proof}
\emph{\olref{equivwo} $\Rightarrow$ \olref{equivcompare}.} Fix $A$ and
$B$. Invoking \olref{equivwo}, there are well-orderings $\tuple{A, R}$
and $\tuple{B, S}$. Invoking
\olref[ordinals][ordtype]{thmOrdinalRepresentation}, let $f \colon
\alpha \to \tuple{A, R}$ and $g \colon \beta \to \tuple{B, S}$ be
isomorphisms. By \olref[sth][ordinals][basic]{ordinalsaresubsets}, either $\alpha \subseteq \beta$ or $\beta \subseteq \alpha$. If $\alpha \subseteq \beta$, then $\comp{f^{-1}}{g} \colon A \to B$ is !!a{injection}, and hence
$\cardle{A}{B}$; similarly, if $\beta \subseteq \alpha$ then $\cardle{B}{A}$.

\emph{\olref{equivcompare} $\Rightarrow$ \olref{equivwo}.} Fix $A$; by
\olref{HartogsLemma} there is some ordinal $\beta$ such that
$\cardnless{\beta}{A}$. Invoking \olref{equivcompare}, we have
$\cardle{A}{\beta}$. So there is some !!{injection} $f \colon A \to
\beta$, and we can use this injection to well-order the elements of
$A$, by defining an order $\Setabs{\tuple{a, b} \in A \times A}{f(a)
\in f(b)}$.
\end{proof}
\noindent
As an immediate consequence: if Well-Ordering fails, then some sets
are \emph{literally incomparable} with regard to their size. So, if
Well-Ordering fails, then transfinite cardinal arithmetic will be
messy. For example, we will have to abandon the idea that if $A$ and
$B$ are infinite then $\cardeq{\cardeq{A \disjointsum B}{A \times
B}}{M}$, where $M$ is the larger of $A$ and $B$ (see
\olref[card-arithmetic][simp]{cardplustimesmax}). The problem is
simple: if we cannot \emph{compare} the size of  $A$ and $B$, then it
is nonsensical to ask which is larger.

% ==== OLP file: content/set-theory/choice/introduction.tex ====


\olfileid{sth}{choice}{intro}

\olsection{Introduction}

In \crefrange{sth:cardinals::chap}{sth:card-arithmetic::chap}, we
developed a theory of cardinals by treating cardinals as ordinals.
That approach depends upon the Axiom of Well-Ordering. It turns out
that Well-Ordering is equivalent to another principle---the Axiom of
Choice---and there has been serious philosophical discussion of its
acceptability. Our question for this chapter are: How is the Axiom
used, and can it be justified?

% ==== OLP file: content/set-theory/choice/justifications.tex ====


\olfileid{sth}{choice}{justifications}
\olsection{Intrinsic Considerations about Choice}

The broader question, then, is whether Well-Ordering, or Choice, or
indeed the comparability of all sets as regards their size---it
doesn't matter which---can be justified. 

Here is an attempted \emph{intrinsic} justification. Back in
\olref[z][story]{sec}, we introduced several principles
about the hierarchy. One of these is worth restating:
\begin{enumerate}
	\item[] \stagesacc. For any stage $S$, and for any sets which were
	formed \emph{before} stage $S$: a set is formed at stage $S$ whose
	members are exactly those sets. Nothing else is formed at
	stage~$S$. 
\end{enumerate}
In fact, many authors have suggested that the Axiom of Choice can be
justified via (something like) this principle. We will briefly provide
a gloss on that approach.

We will start with a simple little result, which offers \emph{yet
another} equivalent for Choice:

\begin{thm}[in $\ZF$]\ollabel{choiceset}
Choice is equivalent to the following principle. If the !!{element}s
of $A$ are disjoint and non-empty, then there is some $C$ such that $C
\cap x$ is a singleton for every $x \in A$. (We call such a $C$ a
{choice set} for $A$.)
\end{thm}

The proof of this result is straightforward, and we leave it as an
exercise for the reader. 

\begin{prob}
Prove \olref[sth][choice][justifications]{choiceset}. If you struggle,
you can find a proof in \cite[pp.~242--3]{Potter2004}.
\end{prob}

The essential point is that a choice set for $A$ is just the range of
a choice function for $A$. So, to justify Choice, we can simply try to
justify its equivalent formulation, in terms of the existence of
choice sets. And we will now try to do exactly that. 

Let $A$'s !!{element}s be disjoint and non-empty. By \stageshier{}
(see \olref[z][story]{sec}), $A$ is formed at some stage~$S$. Note
that all the !!{element}s of $\bigcup A$ are available before stage
$S$. Now, by \stagesacc{}, for \emph{any} sets which were formed
before~$S$, a set is formed whose members are exactly those sets.
Otherwise put: every \emph{possible} collections of earlier-available
sets will exist at~$S$. But it is certainly \emph{possible} to select
objects which could be formed into a choice set for~$A$; that is just
some very specific subset of $\bigcup A$. So: some such choice set
exists, as required.

Well, that's a \emph{very} quick attempt to offer a justification of
Choice on intrinsic grounds. But, to pursue this idea further, you
should read Potter's (\citeyear[\S14.8]{Potter2004}) neat development
of it.

% ==== OLP file: content/set-theory/choice/tarskiscott.tex ====


\olfileid{sth}{choice}{tarskiscott}
\olsection{The Tarski--Scott Trick}

In \olref[cardinals][cardsasords]{defcardinalasordinal}, we
defined cardinals as ordinals. To do this, we assumed the Axiom of
Well-Ordering. We did this, for no other reason than that it is the
``industry standard''.

Before we discuss any of the philosophical issues surrounding
Well-Ordering, then, it is important to be clear that we \emph{can}
depart from the industry standard, and develop a theory of cardinals
\emph{without} assuming Well-Ordering. We can still employ the
definitions of $\cardeq{A}{B}$, $\cardle{A}{B}$ and $\cardless{A}{B}$,
as they appeared in \olref[sfr][siz][]{chap}. We will just need a new
notion of \emph{cardinal}.

A na\"ive thought would be to attempt to define $A$'s cardinality thus:
\[
	\Setabs{x}{\cardeq{A}{x}}.
\]
You might want to compare this with Frege's definition of $\# x Fx$,
sketched at the very end of \olref[cardinals][hp]{sec}. And, for
reasons we gestured at there, this definition fails. Any singleton set
is equinumerous with $\{\emptyset\}$. But new singleton sets are
formed at every successor stage of the hierarchy (just consider the
singleton of the previous stage). So $\Setabs{x}{\cardeq{A}{x}}$ does
not exist, since it cannot have a rank.

To get around this problem, we use a trick due to Tarski and Scott:\footnote{A reminder: all formulas may have parameters (unless explicitly stated otherwise).}

\begin{defn}[Tarski--Scott]
For any formula $\phi(x)$, let
$[ x : \phi(x)] $ be the set of all $x$, of least possible rank, such
that $\phi(x)$ (or $\emptyset$, if there are no $\phi$s).
\end{defn}

We should check that this definition is legitimate. Working in $\ZF$,
\olref[spine][foundation]{zfentailsregularity} guarantees that
$\setrank{x}$ exists for every $x$. Now, if there are any entities
satisfying $\phi$, then we can let $\alpha$ be the least rank such
that  $(\exists x\subseteq V_\alpha)\phi(x)$, i.e., $(\forall \beta
\in \alpha)(\forall x \subseteq V_\beta)\lnot \phi(x)$. We can then
define $[x : \phi(x)]$ by Separation as $\Setabs{x \in
V_{\alpha+1}}{\phi(x)}$. 

Having justified the Tarski--Scott trick, we can now use it to define
a notion of cardinality:

\begin{defn}
The \textsc{ts}-cardinality of $A$ is $\text{tsc}(A) = [x :
\cardeq{A}{x}]$.
\end{defn}

The definition of a \textsc{ts}-cardinal does not use Well-Ordering.
But, even without that Axiom, we can show that
\emph{\textsc{ts}-cardinals} behave rather like \emph{cardinals} as
defined in \olref[cardinals][cardsasords]{defcardinalasordinal}.
For example, if we restate
\olref[cardinals][cardsasords]{lem:CardinalsBehaveRight} and
\olref[card-arithmetic][opps]{lem:SizePowerset2Exp} in terms of
\textsc{ts}-cardinals, the proofs go through just fine in $\ZF$,
without assuming Well-Ordering. 

Whilst we are on the topic, it is worth noting that we can also
develop a theory of ordinals using the Tarski--Scott trick. Where
$\tuple{A, <}$ is a well-ordering, let $\text{tso}(A, <) = [\tuple{X,
R} : \ordeq{\tuple{A, <}}{\tuple{X, R}}]$. For more on this treatment
of cardinals and ordinals, see \citet[chs.~9--12]{Potter2004}.

% ==== OLP file: content/set-theory/choice/vitali.tex ====


\olfileid{sth}{choice}{vitali}
\olsection{Appendix: Vitali's Paradox}

To get a real sense of whether the Banach-Tarski construction is
acceptable or not, we should examine its \emph{proof}. Unfortunately,
that would require much more algebra than we can present here.
However, we can offer some quick remarks which might shed some insight
on the proof of Banach-Tarski,\footnote{For a much fuller treatment,
see \cite{Weston2003} or \cite{Wagon2016}.} by focussing on the
following result:

\begin{thm}[Vitali's Paradox (in $\ZFC$)]\ollabel{vitaliparadox}
Any circle can be decomposed into countably many pieces, which can be
reassembled (by rotation and transportation) to form two copies of
that circle.
\end{thm}

Vitali's Paradox is much easier to prove than the Banach--Tarski Paradox. We have
called it ``Vitali's Paradox'', since it follows from Vitali's
\citeyear{Vitali1905} construction of an unmeasurable set.  But the
set-theoretic aspects of the proof of Vitali's Paradox and the
Banach-Tarski Paradox are very similar. The essential difference
between the results is just that Banach-Tarski considers a
\emph{finite} decomposition, whereas Vitali's Paradox considers a
\emph{countably infinite} decomposition.  As \citet{Weston2003}
puts it, Vitali's Paradox ``is certainly not nearly as striking as the
Banach--Tarski paradox, but it does illustrate that geometric
paradoxes can  happen even in `simple' situations.'' 

Vitali's Paradox concerns a two-dimensional figure, a circle. So we
will work on the plane, $\Real^2$. Let $\rotationsgroup$ be the set of
(clockwise) rotations of points around the origin by \emph{rational}
radian values between $[0,2\pi)$. Here are some algebraic facts about
$\rotationsgroup$ (if you don't understand the statement of the
result, the proof will make its meaning clear):

\begin{lem}\ollabel{rotationsgroupabelian}
$\rotationsgroup$ forms an abelian {group} under composition of functions.
\end{lem}

\begin{proof}
Writing $0_{\rotationsgroup}$ for the rotation by $0$ radians, this is
an identity element for $\rotationsgroup$, since
$\comp{0_{\rotationsgroup}}{\rho} = \comp{\rho}{0_{\rotationsgroup}} =
\rho$ for any $\rho \in \rotationsgroup$.

Every element has an inverse. Where $\rho \in \rotationsgroup$ rotates
by $r$ radians, $\rho^{-1} \in \rotationsgroup$ rotates by $2\pi - r$
radians, so that $\rho \circ \rho^{-1} = 0_\rotationsgroup$.

Composition is associative: $\comp{\rho}{(\comp{\sigma}{\tau})} =
\comp{(\comp{\rho}{\sigma})}{\tau}$ for any $\rho, \sigma, \tau \in
\rotationsgroup$

Composition is commutative: $\comp{\rho}{\sigma} =
\comp{\sigma}{\rho}$ for any $\rho, \sigma \in \rotationsgroup$.
\end{proof}

In fact, we can split our group $\rotationsgroup$
in half, and then use either half to recover the whole group:

\begin{lem}\ollabel{disjointgroup}
There is a partition of $\rotationsgroup$ into two disjoint sets,
$\rotationsgroup_{1}$ and $\rotationsgroup_{2}$, both of which are a
basis for $\rotationsgroup$. 
\end{lem}

\begin{proof}
Let $\rotationsgroup_{1}$ consist of the rotations by rational radian
values in $[0, \pi)$; let $\rotationsgroup_{2} = \rotationsgroup
\setminus \rotationsgroup_1$. By elementary algebra,
$\Setabs{\comp{\rho}{\rho}}{\rho \in \rotationsgroup_1} =
\rotationsgroup$. A similar result can be obtained for
$\rotationsgroup_2$.
\end{proof}

We will use this fact about groups to establish \olref{vitaliparadox}.
Let $\onesphere$ be the unit circle, i.e., the set of points
exactly~$1$ unit away from the origin of the plane, i.e.,
$\Setabs{\tuple{r,s} \in \Real^2}{\sqrt{r^2+s^2}=1}$. We will split
$\onesphere$ into parts by considering the following relation on
$\onesphere$:
\[
	r \sim s \emph{ iff }(\exists \rho \in \rotationsgroup)\rho(r) = s.
\]
That is, the points of $\onesphere$ are linked by this relation iff you can get from one to the other by a rational-valued rotation about the origin. Unsurprisingly:

\begin{lem}
$\sim$ is an equivalence relation.
\end{lem}

\begin{proof}
Trivial, using \olref{rotationsgroupabelian}. 
\end{proof}

We now invoke Choice to obtain a set, $C$, containing exactly one
member from each equivalence class of $\onesphere$ under $\sim$. That
is, we consider a choice function $f$ on the set of equivalence
classes,\footnote{Since $\rotationsgroup$ is !!{enumerable},
each !!{element} of $E$ is !!{enumerable}. Since $\onesphere$ is
!!{nonenumerable}, it follows from \olref{vitalicover} and
\olref[card-arithmetic][simp]{kappaunionkappasize} that $E$ is
!!{nonenumerable}. So this is a use of \emph{un}countable Choice.}
\[
	E = \Setabs{\equivrep{r}{\sim}}{r \in \onesphere},
\]
and let $C = \ran{f}$. For each rotation $\rho \in \rotationsgroup$,
the set $\funimage{\rho}{C}$ consists of the points obtained by
applying the rotation $\rho$ to each point in $C$. These next two
results show that these sets cover the circle completely and without
overlap:

\begin{lem}\ollabel{vitalicover}
$\onesphere = \bigcup_{\rho \in \rotationsgroup} \funimage{\rho}{C}$.
\end{lem}

\begin{proof}
Fix $s \in \onesphere$; there is some $r \in C$ such that $r \in
\equivrep{s}{\sim}$, i.e., $r \sim s$, i.e., $\rho(r) = s$ for some
$\rho \in \rotationsgroup$. 
\end{proof}

\begin{lem}\ollabel{vitalinooverlap}
If $\rho_1 \neq \rho_2$ then $\funimage{\rho_1}{C} \cap \funimage{\rho_2}{C} = \emptyset$. 
\end{lem}

\begin{proof}
Suppose $s \in \funimage{\rho_{1}}{C} \cap \funimage{\rho_{2}}{C}$. So
$s = \rho_{1}(r_{1}) = \rho_{2}(r_{2})$ for some $r_{1}, r_{2} \in C$.
Hence $\rho^{-1}_2(\rho_1(r_1)) = r_2$, and
$\comp{\rho_1}{\rho^{-1}_2} \in \rotationsgroup$, so $r_{1} \sim
r_{2}$. So $r_1 = r_2$, as $C$ selects exactly one member from each
equivalence class under $\sim$. So $s = \rho_1(r_1) = \rho_2(r_1)$,
and hence $\rho_1 = \rho_2$.
\end{proof}

We now apply our earlier algebraic facts to our circle:

\begin{lem}\ollabel{pseudobanachtarski}
There is a partition of $\onesphere$ into two disjoint sets, $D_{1}$
and $D_{2}$, such that $D_{1}$ can be partitioned into countably many
sets which can be rotated to form a copy of $\onesphere$ (and
similarly for $D_{2}$).
\end{lem}

\begin{proof}
Using $\rotationsgroup_{1}$ and $\rotationsgroup_{2}$ from \olref{disjointgroup}, let:
\begin{align*}
	D_{1} &= \bigcup_{\rho \in \rotationsgroup_1} \funimage{\rho}{C} & 
	D_{2} &= \bigcup_{\rho \in \rotationsgroup_2} \funimage{\rho}{C}
\end{align*}
This is a partition of $\onesphere$, by \olref{vitalicover}, and $D_1$
and $D_2$ are disjoint by \olref{vitalinooverlap}. By construction,
$D_1$ can be partitioned into countably many sets,
$\funimage{\rho}{C}$ for each $\rho \in R_1$. And these can be rotated
to form a copy of $\onesphere$, since $\onesphere = \bigcup_{\rho \in
\rotationsgroup}\funimage{\rho}{C} = \bigcup_{\rho \in
\rotationsgroup_1}\funimage{(\comp{\rho}{\rho})}{C}$ by
\olref{disjointgroup} and \olref{vitalicover}. The same reasoning
applies to $D_2$. \end{proof}\noindent This immediately entails
Vitali's Paradox. For we can generate \emph{two} copies of
$\onesphere$ from $\onesphere$, just by splitting it up into countably
many pieces (the various $\funimage{\rho}{C}$'s) and then rigidly
moving them (simply rotate each piece of $D_1$, and first transport
and then rotate each piece of $D_2$).

Let's recap the proof-strategy. We started with some algebraic facts
about the group of rotations on the plane. We used this group to
partition $\onesphere$ into equivalence classes. We then arrived at a
``paradox'', by using Choice to select elements from each class.

We use exactly the same strategy to prove Banach--Tarski. The main
difference is that the algebraic facts used to prove Banach--Tarski
are significantly more complicated than those used to prove
Vitali's Paradox. But those algebraic facts have nothing to do with
Choice. We will summarise them quickly. 

To prove Banach--Tarski, we start by establishing an analogue of
\olref{disjointgroup}: any \emph{free group} can be split into four
pieces, which intuitively we can ``move around'' to recover two copies
of the whole group.\footnote{The fact that we can use \emph{four}
pieces is due to \cite{Robinson1947}. For a recent proof, see
\citet[Theorem 5.2]{Wagon2016}. We follow \citet[p.~3]{Weston2003}
in describing this as ``moving'' the pieces of the group.} We then
show that we can use two particular rotations around the origin of
$\Real^3$ to generate a free group of rotations, $F$.\footnote{See
\citet[Theorem 2.1]{Wagon2016}.} (No Choice yet.) We now regard points
on the surface of the sphere as ``similar'' iff one can be obtained
from the other by a rotation in~$F$. We then \emph{use Choice} to
select exactly one point from each equivalence class of ``similar''
points. Applying our division of $F$ to the surface of the sphere, as
in \olref{pseudobanachtarski}, we split that surface into four pieces,
which we can ``move around'' to obtain two copies of the surface of
the sphere. And this establishes \citep{Hausdorff1914}:

\begin{thm}[Hausdorff's Paradox (in $\ZFC$)] 
The surface of any sphere can be decomposed into finitely many pieces,
which can be reassembled (by rotation and transportation) to form two
disjoint copies of that sphere.
\end{thm}

A couple of further algebraic tricks are needed to obtain the full
Banach-Tarski Theorem (which concerns not just the sphere's surface,
but its interior too). Frankly, however, this is just icing on the
algebraic cake. Hence Weston writes:
\begin{quote}	
  [\ldots] the result on free groups is the \emph{key step} in the
  proof of the Banach-Tarski paradox. From this point of view, the
  Banach-Tarski paradox is not a statement about $\Real^3$ so much as
  it is a statement about the complexity of the group [of translations
  and rotations in $\Real^3$]. \cite[p.~16]{Weston2003}
\end{quote}
That is: whether we can offer a \emph{finite} decomposition (as in
Banach--Tarski) or a \emph{countably infinite} decomposition (as in
Vitali's Paradox) comes down to certain group-theoretic facts about
working in two-dimension or three-dimensions.

Admittedly, this last observation slightly spoils the joke at the end
of \olref[banach]{sec}. Since it is  two dimensional,
``Banach-Tarski'' must be divided into a countable \emph{infinity} of
pieces, if one wants to rearrange those pieces to form ``Banach-Tarski
Banach-Tarski''. To repair the joke, one must write in three
dimensions. We leave this as an exercise for the reader.

One final comment. In \olref[banach]{sec}, we mentioned that the
``pieces'' of the sphere one obtains cannot be \emph{measurable}, but
must be unpicturable ``infinite scatterings''. The same is true of our
use of Choice in obtaining \olref{pseudobanachtarski}. And this is all
worth explaining.

Again, we must sketch some background (but this is \emph{just} a
sketch; you may want to consult a textbook entry on \emph{measure}).
To define a measure for a set $X$ is to assign a value $\mu(E) \in
\Real$ for each $E$ in some ``$\sigma$-algebra'' on $X$. Details here
are not essential, except that the function $\mu$ must obey the
principle of countable additivity: the measure of a countable union of
disjoint sets is the sum of their individual measures, i.e.,
$\mu(\bigcup_{n < \omega} X_n) = \sum_{n < \omega}\mu(X_n)$ whenever
the $X_n$s are disjoint. To say that a set is ``unmeasurable'' is to
say that no measure can be suitably assigned. Now, using our
$\rotationsgroup$ from before:

\begin{cor}[Vitali]
Let $\mu$ be a measure such that $\mu(\onesphere) = 1$, and such that
$\mu(X) = \mu(Y)$ if $X$ and $Y$ are congruent. Then
$\funimage{\rho}{C}$ is unmeasurable for all $\rho \in
\rotationsgroup$. 
\end{cor}

\begin{proof}
For reductio, suppose otherwise. So let $\mu(\funimage{\sigma}{C}) =
r$ for some $\sigma \in \rotationsgroup$ and some $r \in \Real$. For
any $\rho \in C$, $\funimage{\rho}{C}$ and $\funimage{\sigma}{C}$ are
congruent, and hence  $\mu(\funimage{\rho}{C}) = r$ for any $\rho \in
C$. By \olref{vitalicover} and \olref{vitalinooverlap}, $\onesphere =
\bigcup_{\rho \in \rotationsgroup}\funimage{\rho}{C}$ is a countable
union of pairwise disjoint sets. So countable additivity dictates that
$\mu(\onesphere) = 1$ is the sum of the measures of each
$\funimage{\rho}{C}$, i.e.,
\[
	1 = \mu(\onesphere) = \sum_{\rho \in \rotationsgroup}\mu(\funimage{\rho}{C}) = \sum_{\rho \in \rotationsgroup}r
\]
But if $r = 0$ then $\sum_{\rho \in \rotationsgroup}r = 0$, and if $r
> 0$ then $\sum_{\rho \in \rotationsgroup}r = \infty$. 
\end{proof}

% ==== OLP file: content/set-theory/choice/wellorderingproblem.tex ====


\olfileid{sth}{choice}{woproblem}
\olsection{The Well-Ordering Problem}

Evidently rather a lot hangs on whether we accept Well-Ordering. But
the discussion of this principle has tended to focus on an equivalent
principle, the Axiom of Choice. So we will now turn our attention to
that (and prove the equivalence). 

In \citeyear{Cantor1883}, Cantor expressed his support for the Axiom
of Well-Ordering, calling it ``a law of thought which appears to me to
be fundamental, rich in its consequences, and particularly remarkable
for its general validity'' (cited in \citeauthor{Potter2004}
\citeyear[p.~243]{Potter2004}). But Cantor ultimately became convinced
that the ``Axiom'' was in need of proof. So did the mathematical
community. 

The problem was ``solved'' by Zermelo in \citeyear{Zermelo1904}. To
explain his solution, we need some definitions. 

\begin{defn}
A function $f$ is a \emph{choice function} iff $f(x) \in x$ for all $x \in \dom{f}$. We say that $f$ is a \emph{choice function for $A$} iff $f$ is a choice function with $\dom{f} = A \setminus \{\emptyset\}$.
\end{defn}

Intuitively, for every (non-empty) set $x \in A$, a choice function
for $A$ \emph{chooses} a particular element, $f(x)$, from $x$. The
Axiom of Choice is then:

\begin{axiom}[Choice]
	Every set has a choice function.
\end{axiom}

Zermelo showed that Choice entails well-ordering, and vice versa:

\begin{thm}[in $\ZF$]\ollabel{thmwochoice}
Well-Ordering and Choice are equivalent.
\end{thm}

\begin{proof}
\emph{Left-to-right.} Let $A$ be a set of sets. Then $\bigcup A$
exists by the Axiom of Union, and so by Well-Ordering there is some
$<$ which well-orders $\bigcup A$. Now let $f(x) = \text{the $<$-least
member of }x$. This is a choice function for $A$.

\emph{Right-to-left.} Fix $A$. By Choice, there is a choice function,
$f$,  for $\Pow{A} \setminus \{\emptyset\}$. Using Transfinite
Recursion, define a function:
\begin{align*}
	g(0) &= f(A)\\
	g(\alpha) &= 
		\begin{cases}
			\text{stop!{}} &\text{if }A = \funimage{g}{\alpha}\\
			f(A \setminus \funimage{g}{\alpha}) & \text{otherwise}\\	
		\end{cases}
\end{align*}
The indication to ``stop!'' is just a shorthand for what would
otherwise be a more long-winded definition. That is, when $A =
\funimage{g}{\alpha}$ for the first time, let $g(\delta) = A$ for all
$\delta \leq \alpha$. Now, in the first instance, we can only be sure that this defines a \emph{term} (see the remarks after \olref[sth][spine][recursion]{transrecursionschema}); but we will show that we indeed have a function.

Since $f$ is a choice function, for each $\alpha$ (when defined) we have $g(\alpha) =
f(A \setminus \funimage{g}{\alpha}) \in A \setminus
\funimage{g}{\alpha}$; i.e., $g(\alpha) \notin \funimage{g}{\alpha}$.
So if $g(\alpha) = g(\beta)$ then $g(\beta) \notin
\funimage{g}{\alpha}$, i.e., $\beta \notin \alpha$, and similarly
$\alpha \notin \beta$. So $\alpha = \beta$, by Trichotomy. So $g$ is
!!{injective}.

Next, observe that we do stop!{}, i.e.\ that there is some (least) ordinal $\alpha$ such that $A = g[\alpha]$. For suppose otherwise; then as $g$ is !!{injective} we would have $\cardless{\alpha}{\Pow{A} \setminus \{\emptyset\}}$ for
every ordinal $\alpha$, contradicting \olref[hartogs]{HartogsLemma}. Hence also $\ran{g} = A$.

Assembling these facts, $g$ is !!a{bijection} from some ordinal to $A$. Now $g$ can be used to well-order $A$.
\end{proof}

So Well-Ordering and Choice stand or fall together. But the question
remains: do they stand or fall?

